\chapter{The ReasonGuard Framework}
\label{ch:framework}

\section{Three-Layer Architecture}
\label{sec:architecture}

\textsc{ReasonGuard} is structured as a three-layer pipeline. The
\emph{formal layer} emits a verified verdict about a network event in the
form of a JSON formal bound (Section~\ref{sec:bound-schema}). The
\emph{AI layer} consumes the bound through a strictly constrained prompt and
produces a natural-language explanation. The \emph{verification layer}
compares the explanation against the bound and classifies any crossing of
the bound into a reasoning violation class (Section~\ref{sec:taxonomy}).

The three layers are deliberately decoupled. The formal layer can be
replaced with any detector that produces the bound schema, the AI layer can
be any LLM that accepts a system and a user prompt, and the verification
layer can be inspected, audited, and improved independently of either of the
other two. The architecture is designed for the safety-critical setting in
which the formal layer is the source of truth and the AI layer is a
human-readable narrator.

\section{Formal Bound Schema}
\label{sec:bound-schema}

A formal bound is a JSON document with a fixed top-level structure. The
fields are partitioned into five groups: identifying metadata, observable
features, allowed atomic claims, required claims, forbidden claims, and
machine-checkable constraints. Listing~\ref{lst:bound-example} shows a
canonical example produced by \texttt{src/formal\_bound\_builder.py} for one
atomic IEC-104 row.

\begin{lstlisting}[language=Python,caption={Canonical formal bound (proxy
schema v2.1).},label={lst:bound-example}]
{
  "event_id": "evt_000001",
  "dataset": "eon_iec",
  "event_granularity": "atomic_row",
  "timestamp_ms": 1654646400041,
  "formal_bound_version": "2.1",
  "schema_status": "proxy_until_official_automaton_schema",
  "protocol_hint": "iec104",
  "communication_direction": "client_to_server",
  "frame_type": "s_or_u_format_control",
  "observation_pattern": "iec104_control_frame_pattern",
  "formal_class": "IEC104_CONTROL_FRAME_OBSERVATION",
  "formal_class_status": "proxy_not_ground_truth_attack_label",
  "features": {
    "source_id": 1, "destination_id": 2, "bytes": 46, "pkt_length": 4,
    "srcport": 57528, "dstport": 2404,
    "asdu_address": 0, "asdu_cot": 0,
    "asdu_items": 0, "asdu_type": 0, "frame_fmt": 1
  },
  "allowed_facts": { ... },
  "allowed_claims": [ ten atomic statements ],
  "required_claims": [ five must-mention statements ],
  "forbidden_claims": [ eight must-not-mention statements ],
  "machine_constraints": {
    "must_match": { ... per-feature equality requirements },
    "attack_claim_allowed": false,
    "malicious_claim_allowed": false,
    "mitigation_claim_allowed": false,
    "causal_claim_allowed": false,
    "normality_claim_allowed": false,
    "scope_expansion_allowed": false
  },
  "proxy_event_type": "topology_change_proxy_until_official_schema",
  "thesis_note": "..."
}
\end{lstlisting}

The schema is intentionally conservative under the proxy regime: every flag
in \texttt{machine\_constraints} defaults to \texttt{false} until the
official schema explicitly permits it. This guarantees that the proxy bound
can only mark an explanation as violating; it cannot falsely mark an
explanation as fully cleared. Once the official NES@FIT automaton schema is
integrated, the corresponding flags will be lifted to \texttt{true} where
appropriate (for example, when the formal verdict establishes a confirmed
classification), and the bound versioning in \texttt{formal\_bound\_version}
will be incremented to track the schema revision.

\subsection{Proxy Event Types}
\label{subsec:proxy-event-types}

The exposé identifies five canonical ICS event types: normal Modbus polling,
function code violation, network scan, replay attack, and topology change.
Until the official automaton output is available, these labels are assigned
by the proxy classifier in \texttt{src/event\_type\_classifier.py}, which
applies a small rule set to the observable features. Every proxy label is
suffixed with \texttt{\_proxy\_until\_official\_schema} so the downstream
artefacts retain unambiguous provenance.

\section{Reasoning Violation Taxonomy}
\label{sec:taxonomy}

The taxonomy is the primary theoretical contribution of this thesis. Five
violation classes are defined with respect to the bound:

\begin{description}
  \item[V1 Fabricated reasoning.] The AI asserts facts, causes, or entities
        that are not present in or derivable from the formal output. The
        archetypal failure is the inference of an attack class or a causal
        explanation that the bound does not support.
  \item[V2 Contradicted reasoning.] The AI explicitly contradicts a fact
        established by the bound. The most common case in practice is the
        inversion of the communication direction
        (\texttt{server\_to\_client} where the bound states
        \texttt{client\_to\_server}, or vice versa).
  \item[V3 Over-generalised reasoning.] The AI produces a correct
        high-level category but loses formally established specifics that
        the operator needs in order to act. The archetypal failure is the
        replacement of specific numeric identifiers with descriptive
        adjectives.
  \item[V4 Under-specified reasoning.] The AI omits formally established
        information that is operationally required for correct response.
        Unlike V3 the explanation does not over-generalise; it simply leaves
        the formally available evidence unused.
  \item[V5 Incoherent reasoning.] The AI explanation is internally
        self-contradictory, independent of the formal output. The
        archetypal failure is the simultaneous claim that the same event is
        normal and malicious.
\end{description}

\subsection{Severity Classification}
\label{subsec:severity}

Each violation set is mapped to one of four severity levels:

\begin{itemize}
  \item \textbf{High.} V1, V2, or V5 fires. The explanation is either
        unsupported, contradicting the bound, or internally incoherent.
  \item \textbf{Medium.} V3 and V4 fire together. The explanation is
        actionable in category terms but is missing operational detail.
  \item \textbf{Low.} V3 or V4 fires alone.
  \item \textbf{None.} No violation fires.
\end{itemize}

\section{Verification Mechanism}
\label{sec:detection}

For each (formal bound, AI explanation) pair, \textsc{ReasonGuard} runs a
three-step verification pipeline.

\subsection{Structured Claim Extraction}
\label{subsec:claim-extraction}

The first step uses spaCy with a custom \texttt{EntityRuler} for the ICS
domain (\texttt{src/claim\_extraction/entity\_ruler.py}). The ruler defines
patterns for the protocol (\texttt{PROTOCOL}), the communication direction
(\texttt{DIRECTION}), the IEC-104 frame type (\texttt{FRAME\_TYPE}), ASDU
fields (\texttt{ASDU\_FIELD}), network endpoints
(\texttt{NETWORK\_ENDPOINT}), attack categories (\texttt{ATTACK\_CATEGORY}),
mitigation terms (\texttt{MITIGATION}), normality terms
(\texttt{NORMALITY}), and uncertainty markers (\texttt{UNCERTAINTY}). The
ruler is registered before the statistical NER component so the deterministic
ICS patterns take precedence on the relevant spans, while the statistical
model continues to label generic entities such as numeric values.

\subsection{Bound Comparison}
\label{subsec:bound-comparison}

The second step compares each extracted claim against the bound and assigns
it one of three verdicts:

\begin{description}
  \item[SUPPORTED.] The claim is consistent with the bound.
  \item[UNSUPPORTED.] The claim is not consistent with the bound and is not
        explicitly contradicted by it.
  \item[CONTRADICTED.] The claim is explicitly negated by the bound.
\end{description}

The per-claim verdict is implemented in
\texttt{src/claim\_extraction/three\_state\_classifier.py}. It checks the
direction terms, the source / destination / port / ASDU references, the
attack / mitigation / normality term sets against the
\texttt{machine\_constraints} flags, and the numeric identifiers against
the union of bound-derived numbers and a small set of protocol constants.

\subsection{Violation Classification}
\label{subsec:violation-classification}

The third step maps the pattern of per-claim verdicts to the V1--V5
taxonomy using a rule-based classifier
(\texttt{src/reason\_guard\_checker.py}). The classifier has two layers:

\begin{enumerate}
  \item A lexicon layer that uses curated term lists for the attack,
        mitigation, causal, and normality categories. The lexicon layer is
        always active and provides the conservative baseline.
  \item A spaCy dependency-parse layer that uses the parse tree to
        distinguish between assertion and negation. A forbidden term is
        treated as safely negated when the parser can connect it to a
        negation token within the same sentence, or when the sentence
        contains a strong descriptive marker (\texttt{proxy},
        \texttt{schema}).
\end{enumerate}

The two layers are composed: the spaCy layer filters the candidate term
list before the lexicon assertions are evaluated. When spaCy is not
installed the framework falls back to a substring-pattern allowlist
(\texttt{SAFE\_NEGATION\_PATTERNS}); the substring patterns cover the
canonical phrasings observed in the synthetic regression set but not the
multi-clause negations observed on real local-LLM outputs.

\subsection{V5 Sentence-Level Coherence}
\label{subsec:v5-coherence}

V5 (incoherent reasoning) is detected using spaCy's sentence segmentation
(\texttt{src/claim\_extraction/coherence.py}). The detector scans each
sentence for the canonical conflicting pairs (for example,
\emph{normal} and \emph{malicious}, or \emph{client-to-server} and
\emph{server-to-client}) and reports a finding only when both members of a
pair appear within the same sentence and the sentence does not already
contain a negation marker. This avoids the false positives produced by a
naive whole-document substring search.

\section{Grounding in Runtime Verification Theory}
\label{sec:grounding}

The verification mechanism is grounded in the runtime-verification
framework of \autocite{LeuckerSchallhart2009RuntimeVerification}. The
property being verified is reasoning faithfulness: the AI explanation must
not assert facts inconsistent with the formal verdict. The formal automaton
output serves as the runtime monitor specification, and the AI explanation
is the execution under inspection. This grounding distinguishes
\textsc{ReasonGuard} from post-hoc quality scoring approaches: it is a
runtime verification system, not a quality scoring system, and a
\textsc{ReasonGuard} verdict can in principle be gated on at deployment
time rather than only reported offline.

\section{Implementation Notes}
\label{sec:implementation-notes}

The implementation is structured as a Python package under \texttt{src/}.
The pipeline modules are:

\begin{itemize}
  \item \texttt{pipeline\_config.py} — paths, model lists, MLflow names,
        seeds.
  \item \texttt{formal\_bound\_builder.py} — CSV $\to$ JSON bounds.
  \item \texttt{event\_type\_classifier.py} — proxy event-type assignment.
  \item \texttt{prompt\_builder.py} — bounds $\to$ strict prompts (with
        optional stratified sampling).
  \item \texttt{synthetic\_response\_generator.py} — adversarial
        regression set.
  \item \texttt{ollama\_llm\_runner.py} — local-model inference via Ollama.
  \item \texttt{cloud\_llm\_runner.py} — OpenAI and Gemini clients.
  \item \texttt{reason\_guard\_checker.py} — V1--V5 verification.
  \item \texttt{claim\_extraction/} — spaCy NER, three-state classifier,
        negation scope, coherence.
  \item \texttt{evaluation\_metrics.py} — precision, recall, F1, Kappa.
  \item \texttt{event\_analysis.py} — model $\times$ event-type breakdown.
  \item \texttt{viz.py} — matplotlib figures.
  \item \texttt{manual\_annotation.py} — annotation CSV builder.
  \item \texttt{dashboard/} — Streamlit dashboard skeleton.
\end{itemize}

Every module is unit-tested in \texttt{tests/} and every run is tracked in
MLflow under the experiments declared in
\texttt{src/pipeline\_config.py}.
